\chapter{Column physics on GPUs}\label{ch:colphys}

The physics schemes of \cref{ch:mp,ch:rad,ch:sfc,ch:lsm,ch:pbl} work on one column at a time. Their columns are
independent, so they parallelize naturally over columns, but each column is a long sequential computation with many
temporary values per level. This is the opposite of a stencil: little memory traffic per operation, much state per
thread. This chapter looks at what limits such kernels on a GPU, and at the port's two designs: one thread per column
(Template CP) for most schemes, and a batched, restructured RRTMG for the longwave radiation.

\section{One thread per column}

Template CP (\cref{ch:port}) gives each $(i,j)$ column to one thread, which copies its column into private arrays,
calls the scheme for a single column, and copies the results back. The parallelism is the number of columns:
\begin{itemize}
  \item 656\,100 on d02 ($810 \times 810$);
  \item 201\,601 on d01 ($449 \times 449$).
\end{itemize}
An A100 holds at most 221\,184 resident threads (108 SMs $\times$ 2048) and an H100 SXM 270\,336. d02 fills the GPU
three times over; d01 about once, if the occupancy were full, which for these kernels it is not.

\section{What limits a column kernel}

\paragraph{State per thread.} WSM6 keeps about 78 arrays of one column level each, plus the sedimentation arrays; Noah
and YSU keep dozens. Most of this lives in local memory (\cref{ch:gpuarch}): per thread in device memory, cached in L1
and L2. Thanks to the interleaving of local memory, accesses at the same level by the threads of a warp are coalesced
(\cref{ch:memory}).

\paragraph{The local-memory reservation.} The GPU reserves local memory for every thread that can be resident at
once, whether it runs or not:
\begin{equation}
  M_{\text{local}} = s_{\text{frame}} \times n_{\text{resident}},
\end{equation}
where $s_{\text{frame}}$ is the largest per-thread stack frame of any kernel, which the compiler reports
(\code{-gpu=ptxinfo}). With about 30~KB per thread for WSM6, that is about 6.6~GB on an A100 and 8.1~GB on an H100
SXM. The memory estimate of the case reserves 8.3~GB for it (\cref{ch:perf}). The device stack limit must be at least
the largest frame; the port sets it from the measured frames (\code{NV_ACC_CUDA_STACKSIZE}) and checks it at startup.

\paragraph{Fixed sizes.} The private arrays have a fixed size, \code{WRF_KMAX} = 64 levels, for 59 used: arrays sized
at run time would live in the device heap, which is slow and small. The 8\,\% of unused levels cost reservation, not
traffic.

\paragraph{Registers and occupancy.} A column kernel uses many registers, often the maximum, so few warps fit on an
SM (\cref{eq:occ-regs}). With few warps, latency is poorly hidden, and each thread's long sequence of dependent
operations runs at the latency of local memory. The best register cap and launch shape differ between the A100 and
the H100 (\cref{ch:warp}).

\paragraph{Divergence and imbalance.} Columns do different amounts of work: cloudy columns run microphysical processes
that clear columns skip; the sedimentation of WSM6 takes as many sub-steps as the fastest-falling hydrometeor needs; a
night column leaves the shortwave scheme at once. A warp runs as long as its slowest column.

\section{RRTMG: a batched design}

The longwave scheme RRTMG \citep{iacono2008} does not fit one thread per column. For each column it computes optical
depths for 140 spectral $g$-points on 109 layers (the model's layers plus the layers it adds above the model top), and
then a radiative transfer for each $g$-point (\cref{ch:rad}). Its temporary storage is about 1.1~MB per column. For all
656\,100 columns of d02 at once, that would be 720~GB.

\paragraph{Correctness version.} The port processes the columns in batches of $B = 4096$ (\code{WRF_RRTMG_BATCH}),
with the per-column temporaries hoisted into batch work arrays with a trailing batch index: 4.5~GB. A kernel runs one
thread per column of the batch, and each thread executes the original code for its column, unchanged apart from its
declarations. It is correct, and slow: 4096 threads use 2\,\% of the threads an A100 can hold, and d02 needs 161 such
batches per radiation call.

\paragraph{Performance version.} The work inside a column has more parallelism than the code shows: the optical depths
of different layers are independent, and the radiative transfer of different $g$-points is independent until their
fluxes are added. Optimization O6 (phase S6-13) splits the column into three kernels (\cref{alg:rrtmg-perf}).

\begin{algorithm}[ht]
\caption{The performance version of RRTMG longwave for a batch of $B$ columns (phase S6-13; class R2).}
\label{alg:rrtmg-perf}
\begin{algorithmic}[1]
\State \textbf{K-RRTMG-TAU}, one thread per (column $b$, layer): the optical depths of all $g$-points (\code{taumol}, \code{taugb1} \dots\ \code{taugb16})
\State \textbf{K-RRTMG-RT1}, one thread per (column $b$, $g$-point): the downward recurrence from the top, the
  surface reflection, the upward recurrence, as in \code{rtrnmc} for that $g$-point; store the radiances
  $d_{\ell,g}$ and $u_{\ell,g}$ of every level $\ell$
\State \textbf{K-RRTMG-RT2}, one thread per (column $b$, level $\ell$):
\For{each band, in order}
  \State $D \gets 0$; \quad $U \gets 0$
  \For{each $g$-point of the band, in order}
    \State $D \gets D + d_{\ell,g}$; \quad $U \gets U + u_{\ell,g}$
  \EndFor
  \State $F^{\downarrow}_\ell \gets F^{\downarrow}_\ell + (D\,w_{\text{diff}})\,\Delta\nu_{\text{band}}$; \quad
    the same for $F^{\uparrow}_\ell$
\EndFor
\end{algorithmic}
\end{algorithm}

Why are the bits unchanged? In \code{rtrnmc}, the radiances of each $g$-point depend only on that $g$-point's optical
depths and Planck fractions, so computing them in separate threads gives the same values. The fluxes are sums: within
a band, the source adds the radiances of the $g$-points one after the other into \code{drad} and \code{urad}; after the
band, it multiplies by the diffusivity factor and adds the result, times the band width, to the total fluxes. Kernel
RT2 performs exactly these additions, in exactly this order. The parallel work computes the terms; the sums are still
sequential. This is class R2 (\cref{tab:bitclass}): parallel work, ordered accumulation.

The performance version has $4096 \times 140 = 573\,440$ threads in RT1, enough to fill either GPU, and its tests are
the column harness T-RRTMG-COL (host against device on $10^5$ columns from the reference state) and the radiation
window T-TRACE-RAD.

\section{WSM6 and Noah}

For the other schemes the questions are of shape, not of structure:
\begin{itemize}
  \item \textbf{Batch shape.} A column kernel over all columns at once needs the full local-memory reservation; a
    kernel over a part of the domain at a time needs less memory but more launches. Phase S6-13.4 measures both for
    WSM6 on each GPU.
  \item \textbf{Splitting.} A scheme can be split at a point where all columns have finished one phase (for WSM6,
    before and after sedimentation), with the intermediate values in work arrays. Fewer live values per kernel mean
    fewer registers and spills (\cref{ch:warp}); the statements of each column keep their order.
  \item \textbf{Level parallelism.} Parts of a scheme that are pointwise in $k$ (the saturation computations of WSM6,
    the slope parameters) could run one thread per level. The recurrences in $k$ (sedimentation, the implicit
    vertical diffusion of YSU) cannot.
\end{itemize}

\begin{keypoint}
Column physics is where a GPU's strengths matter least and the bitwise rule constrains most: the recurrences and sums
in $k$ must stay sequential. The gains come from more parallelism around the recurrences (RRTMG's $g$-points), from
less state per thread, and from the right launch shape for each GPU.
\end{keypoint}
